\begin{lemma}
\label{lemma-lift-to-system-complexes}
Let $(A_n)$ be an inverse system of rings. Suppose that we are given
\begin{enumerate}
\item for every $n$ an object $K_n$ of $D(A_n)$, and
\item for every $n$ a map $\varphi_n : K_{n + 1} \to K_n$ of
$D(A_{n + 1})$ where we think of $K_n$ as an object of $D(A_{n + 1})$
by restriction via $A_{n + 1} \to A_n$.
\end{enumerate}
There exists an object
$M = (M_n^\bullet) \in D(\textit{Mod}(\mathbf{N}, (A_n)))$
and isomorphisms $\psi_n : M_n^\bullet \to K_n$ in $D(A_n)$
such that the diagrams
$$
\xymatrix{
M_{n + 1}^\bullet \ar[d]_{\psi_{n + 1}} \ar[r] & M_n^\bullet \ar[d]^{\psi_n} \\
K_{n + 1} \ar[r]^{\varphi_n} & K_n
}
$$
commute in $D(A_{n + 1})$.
\end{lemma}

\begin{proof}
We write out the proof in detail. For an $A_n$-module $T$ we write
$T_{A_{n + 1}}$ for the same module viewd as an $A_{n + 1}$-module.
Suppose that $K_n^\bullet$ is a complex of $A_n$-modules representing $K_n$.
Then $K_{n, A_{n + 1}}^\bullet$ is the same complex, but viewed as a
complex of $A_{n + 1}$-modules. By the construction of the derived
category, the map $\psi_n$ can be given as
$$
\psi_n = \tau_n \circ \sigma_n^{-1}
$$
where $\sigma_n : L_{n + 1}^\bullet \to K_{n + 1}^\bullet$
is a quasi-isomorphism of complexes of $A_{n + 1}$-modules
and $\tau_n : L_{n + 1}^\bullet \to K_{n, A_{n + 1}}^\bullet$
is a map of complexes of $A_{n + 1}$-modules.

\medskip\noindent
Now we construct the complexes $M_n^\bullet$ by induction. As base case
we let $M_1^\bullet = K_1^\bullet$. Suppose we have already constructed
$M_e^\bullet \to M_{e - 1}^\bullet \to \ldots \to M_1^\bullet$
and maps of complexes $\psi_i : M_i^\bullet \to K_i^\bullet$ such that
the diagrams
$$
\xymatrix{
M_{n + 1}^\bullet \ar[d]_{\psi_{n + 1}} \ar[rr] & &
M_{n, A_{n + 1}}^\bullet \ar[d]^{\psi_{n, A_{n + 1}}} \\
K_{n + 1}^\bullet & L_{n + 1}^\bullet \ar[l]_{\sigma_n} \ar[r]^{\tau_n} &
K_{n, A_{n + 1}}^\bullet
}
$$
above commute in $D(A_{n + 1})$ for all $n < e$. Then we consider the diagram
$$
\xymatrix{
& & M_{e, A_{e + 1}}^\bullet \ar[d]^{\psi_{e, A_{e + 1}}} \\
K_{e + 1}^\bullet & L_{e + 1}^\bullet \ar[r]^{\tau_e} \ar[l]_{\sigma_e} &
K_{e, A_{e + 1}}^\bullet
}
$$
in $D(A_{e + 1})$. Because $\psi_e$ is a quasi-isomorphism,
we see that $\psi_{e, A_{e + 1}}$ is a quasi-isomorphism too.
By the definition of morphisms in $D(A_{e + 1})$ we can
find a quasi-isomorphism
$\psi_{e + 1} : M_{e + 1}^\bullet \to K_{e + 1}^\bullet$
of complexes of $A_{e + 1}$-modules such that there exists a morphism
of complexes $M_{e + 1}^\bullet \to M_{e, A_{e + 1}}^\bullet$
of $A_{e + 1}$-modules representing the composition 
$\psi_{e, A_{e + 1}}^{-1} \circ \tau_e \circ \sigma_e^{-1}$
in $D(A_{e + 1})$. Thus the lemma holds by induction.
\end{proof}